% ===== UNIT 35: Why fusion works + backbone scale (6.2-6.3) =====
% Status: AI-DRAFT  (change to REWRITTEN / REVIEWED when done; see CHECKLIST.md)
% Rules: style rewrite only. Keep every \cite, \ref, \label, number, equation and table value as is. Change the wording, not the content.
\subsection{Why score-level fusion works}
The complementarity is architectural. The ViT-B backbone processes images as
global patch-token sequences with multi-head self-attention, producing
representations sensitive to holistic vehicle structure and long-range spatial
relationships~\cite{he2021transreid,dosovitskiy2021vit}. The CLIP-SENet
ResNet101-IBN-a backbone~\cite{pan2018ibn} applies hierarchical local
convolutional filters capturing fine-grained texture, with IBN providing
simultaneous style normalisation and discriminative training. The AFEM branch
injects attribute-level structure from TinyCLIP~\cite{lu2025clipsenet}
orthogonal to ViT patch attention. Score-level fusion preserves each stream's
post-processing calibration independently and requires no additional training.
The optimal $w_2{=}0.65$ for CLIP-SENet reflects its higher standalone
post-processed performance (91.44\% vs.\ 89.97\% mAP for ViT-B); the ViT-B
branch still contributes $+1.88$\,pp above CLIP-SENet alone, confirming genuine
complementarity.

\subsection{Backbone scale and overfitting}
ViT-L/14 (304\,M parameters) was trained under the same recipe as a
higher-capacity alternative. Despite its larger parameter count it overfit
VeRi-776's 576 training identities---a known failure mode at this dataset
scale---and did not exceed the ViT-B/16 CLIP result after re-ranking. This
supports prioritising systematic post-processing over backbone scaling as the
primary performance lever for small-scale ReID benchmarks.
